\documentclass[10pt,twocolumn]{article}
\usepackage[margin=0.62in,columnsep=0.24in]{geometry}
\usepackage[T1]{fontenc}
\usepackage{lmodern}
\usepackage{amsmath,amssymb,mathtools,bm}
\usepackage{xcolor}
\usepackage{microtype}
\usepackage[colorlinks=true,linkcolor=blue,urlcolor=blue,citecolor=blue]{hyperref}

\setlength{\parindent}{0pt}
\setlength{\parskip}{3pt}
\setlength{\abovedisplayskip}{5pt}
\setlength{\belowdisplayskip}{5pt}
\setlength{\abovedisplayshortskip}{3pt}
\setlength{\belowdisplayshortskip}{3pt}

\newcommand{\blk}[1]{\vspace{3pt}\noindent\textbf{#1:}\quad}
\newcommand{\Ek}{E_k}
\newcommand{\Ok}{\mathcal O_k}
\newcommand{\Edel}{\mathcal E_k^{\ominus}}
\newcommand{\Gdel}{\mathcal G_k^{\ominus}}
\newcommand{\Ndel}{\mathcal N_k^{\ominus}}
\newcommand{\Cdel}{\mathcal C_k^{\ominus}}
\newcommand{\Lam}{\Lambda_j^{\ominus}}
\newcommand{\nudel}{\nu_j^{\ominus}}
\newcommand{\Rdel}{R_j^{\ominus}}
\newcommand{\dEdel}{\Delta E_j^{\ominus}}
\newcommand{\clip}{\Pi_{[0,1]}}

\begin{document}

\title{\vspace{-2em}Paper-I SNAKE Prune Redesign: Pedagogical Math Note}
\author{Planning companion to \texttt{paper\_i\_snake\_prune\_redesign\_plan\_20260704.md}}
\date{2026-07-04}
\maketitle
\vspace{-1.2em}

\textbf{Scope.} This note is a mathematical design aid for the prune redesign. It is not an implementation patch and does not change Paper-I results.

\section*{1. Present Prune Contract}

\blk{Micro-intuition}
The current prune path first asks which existing ansatz coordinates may be considered, then uses Schur only to choose what to measure.

\blk{Mathematical derivation}
At adaptive step \(k\), the active ordered ansatz and energy are
\[
\Ok=(d_1,\ldots,d_{n_k}),
\qquad
\Ek=E(\Ok,\boldsymbol\theta_k).
\]
Here \(k\) is the current adaptive step, \(\Ok\) is the ordered ansatz operator sequence at that step, \(d_j\) is its \(j\)-th admitted generator, \(n_k\) is the current ansatz length, \(\boldsymbol\theta_k\) is the optimized parameter vector on that ansatz, and \(\Ek\) is the corresponding variational energy.
Coordinate \(d_j\) has admission step \(a_j\) and cooldown \(c_j(k)\). The current deletion-eligible set is
\[
\boxed{
\Edel=
\{j\in\{1,\ldots,n_k\}: k-a_j\ge p_{\rm protect},\ c_j(k)\le 0\}.
}
\]
Here \(\Edel\) is the set of coordinate indices that prune is allowed to consider, \(a_j\) is the step when \(d_j\) was first admitted, \(p_{\rm protect}\) is the minimum protected age before deletion can be attempted, and \(c_j(k)\) is the remaining cooldown for coordinate \(j\). The condition \(c_j(k)\le0\) means the cooldown has expired.
For \(j\in\Edel\), the Schur surrogate produces a predicted deletion loss \(\Lam\). The measured nonlinear deletion loss is
\[
\boxed{
\dEdel
=
E(\Ok\ominus d_j,\boldsymbol\theta_{k,-j})-\Ek .
}
\]
Here \(\Ok\ominus d_j\) is the ansatz after deleting generator \(d_j\), \(\boldsymbol\theta_{k,-j}\) is the parameter vector used for that delete-refit trial, and \(\dEdel\) is the actual measured energy regression caused by deletion. Negative or small \(\dEdel\) means the deletion is energetically safe.
The authoritative deletion accept/reject variable has the form
\[
\boxed{
A_j^\ominus
=
\mathbf 1\{\dEdel\le\epsilon_k^\ominus\}
\mathbf 1\{G_j^{\rm ret}\ge\rho_{\rm ret}G_k^{\rm adm}\}
\mathbf 1\{C_j^{\rm curv}=1\}.
}
\]
Here \(A_j^\ominus=1\) means the deletion is accepted, \(\epsilon_k^\ominus\) is the allowed prune energy-regression tolerance, \(G_j^{\rm ret}\) is the post-deletion retained gain, \(G_k^{\rm adm}\) is the gain from the just-accepted admission, \(\rho_{\rm ret}\) is the required retained-gain fraction, and \(C_j^{\rm curv}\) is the curvature-safety indicator.

\blk{Physical interpretation}
Schur controls which candidate is worth paying to test. The measured delete-refit energy test decides whether deletion is actually committed.

\section*{2. Defect in the Current Nomination Key}

\blk{Micro-intuition}
The lexicographic key behaves like a deterministic sorting ladder, not like one coherent prune utility.

\blk{Mathematical derivation}
The current ranking object can be summarized as
\[
K_j=
(\Lam,b_j,\alpha_j,-s_j,\beta_j,|\bar\theta_j|,\ell_j),
\qquad
K_j<K_i\Rightarrow j\succ i.
\]
Here \(K_j\) is the current lexicographic prune ranking key for coordinate \(j\), \(b_j\) is the stored admission/selector score discounted by stored burden, \(\alpha_j\) is the small-angle pool flag, \(s_j\) is the smallness or stagnation score, \(\beta_j\) is the stored selector burden, \(\bar\theta_j\) is the wrapped angle, \(\ell_j\) is the coordinate label, and \(j\succ i\) means \(j\) is ranked ahead of \(i\).
Lexicographic ordering means
\[
\Lam<\Lambda_i^\ominus
\quad\Longrightarrow\quad
j\succ i
\]
Here \(\Lambda_i^\ominus\) is the Schur predicted deletion loss for the comparison coordinate \(i\).
regardless of the other entries. The later entries matter only after earlier entries tie or nearly tie in the sorting implementation.

\blk{Physical interpretation}
A prune decision should trade saved cost, predicted loss, redundancy, history, and conditioning in one scalar pressure. A long tie-breaker stack is not the mathematical object we want to tune.

\section*{3. Replacement: Hard Gates Plus Scalar Score}

\blk{Micro-intuition}
First remove coordinates that are physically or numerically unsafe to nominate; then score the remaining coordinates by one scalar deletion utility.

\blk{Mathematical derivation}
Define the hard-gated nomination domain
\[
\boxed{
\begin{gathered}
\Gdel=
\{j\in\Edel:
C_j^{\rm saved}>0,\ \Lam<\infty,\\
\nudel\le\nu_{\max},\
\Delta\kappa_j\le\Delta\kappa_{\max}\}.
\end{gathered}
}
\]
Here \(\Gdel\) is the subset of eligible coordinates that pass hard prune-safety gates. The index \(j\) denotes a candidate deletion coordinate. \(C_j^{\rm saved}\) is the estimated compiled-resource or measurement burden saved by deleting \(d_j\). \(\Lam\) is the Schur predicted deletion loss. \(\nudel\) is the deletion-novelty residual: large values mean the deleted direction is not well represented by survivors. \(\nu_{\max}\) is the maximum allowed deletion novelty. \(\Delta\kappa_j\) is the estimated conditioning damage from deleting \(d_j\), and \(\Delta\kappa_{\max}\) is the maximum allowed conditioning damage.
For \(j\in\Gdel\), define
\[
\boxed{
\Rdel
=
\frac{
C_j^{\rm saved}
w_j^{\rm small}
w_j^{\rm stale}
w_j^{\rm red}
}{
\epsilon_R
+\Lambda_{j,+}^{\ominus}
+\lambda_\nu\nudel
+\lambda_{\rm rej}\bar L_j^{\rm rej}
+\lambda_\kappa\Delta\kappa_j
},
}
\]
Here \(\Rdel\) is the scalar prune nomination score. The factors \(w_j^{\rm small}\), \(w_j^{\rm stale}\), and \(w_j^{\rm red}\) respectively favor small-angle, old, and redundant coordinates. The denominator penalizes predicted deletion loss \(\Lambda_{j,+}^{\ominus}\), deletion novelty \(\nudel\), rejection history \(\bar L_j^{\rm rej}\), and conditioning damage \(\Delta\kappa_j\). The positive constants \(\epsilon_R,\lambda_\nu,\lambda_{\rm rej},\lambda_\kappa\) are tunable score-scale and penalty weights.
The clipped Schur loss is
\[
\Lambda_{j,+}^{\ominus}:=\max\{0,\Lam\}.
\]
Here \(\Lambda_{j,+}^{\ominus}\) is the nonnegative part of the Schur predicted deletion loss, so a negative surrogate loss is treated as zero risk rather than as a negative denominator contribution.
The unpersisted top nomination is
\[
\boxed{
j_k^\star\in\arg\max_{j\in\Gdel}\Rdel .
}
\]
Here \(j_k^\star\) is the coordinate index with the largest prune nomination score at step \(k\), before applying the persistence gate.

\blk{Physical interpretation}
The numerator says why deletion is attractive. The denominator says why deletion is risky. Negative Schur predictions do not make the score singular because \(\Lambda_{j,+}^{\ominus}\) is clipped at zero.

\section*{4. Saved Burden}

\blk{Micro-intuition}
Deletion should be more attractive when it saves compiled resources or measurement burden.

\blk{Mathematical derivation}
Let deleting \(d_j\) save the nonnegative resource amounts
\[
\Delta N_{2q,j}^+,\qquad
\Delta D_{2q,j}^+,\qquad
\Delta D_{c,j}^+,\qquad
\Delta S_j^+,
\]
where \(x^+=\max\{0,x\}\).
Here \(\Delta N_{2q,j}^+\) is the saved two-qubit gate count, \(\Delta D_{2q,j}^+\) is the saved two-qubit depth, \(\Delta D_{c,j}^+\) is the saved total circuit depth, and \(\Delta S_j^+\) is the saved estimator-work or shot-count proxy from deleting \(d_j\). The superscript \(+\) keeps only positive savings, so a deletion that does not save a resource contributes zero for that resource.
A normalized saved-burden score is
\[
\boxed{
\begin{aligned}
C_j^{\rm saved}
&=
\omega_N\frac{\Delta N_{2q,j}^{+}}{N_0+\epsilon_C}
+\omega_{D2}\frac{\Delta D_{2q,j}^{+}}{D_{2,0}+\epsilon_C}
\\
&\quad
+\omega_{Dc}\frac{\Delta D_{c,j}^{+}}{D_{c,0}+\epsilon_C}
+\omega_S\frac{\Delta S_j^{+}}{S_0+\epsilon_C}.
\end{aligned}
}
\]
Here \(C_j^{\rm saved}\) is a dimensionless weighted sum of saved resources. The constants \(\omega_N,\omega_{D2},\omega_{Dc},\omega_S\) choose the relative importance of the four resources; \(N_0,D_{2,0},D_{c,0},S_0\) are normalization scales; and \(\epsilon_C>0\) prevents division by zero.

\blk{Physical interpretation}
\(C_j^{\rm saved}\) is the resource side of prune. It is allowed to prefer deleting a costly generator, but only if the denominator says the deletion is geometrically and energetically plausible.

\section*{5. Smooth Small-Angle and Stale-Age Pressures}

\blk{Micro-intuition}
Small amplitudes and old coordinates are weak deletion evidence, so they should be smooth weights rather than binary authorities.

\blk{Mathematical derivation}
Use the wrapped coordinate magnitude
\[
\bar\theta_j=\operatorname{wrap}_{[-\pi,\pi)}(\theta_j),
\qquad
a_j^\theta=|\bar\theta_j|.
\]
Here \(\theta_j\) is the current variational angle for coordinate \(j\), \(\bar\theta_j\) is that angle wrapped into the interval \([-\pi,\pi)\), and \(a_j^\theta\) is its absolute wrapped magnitude.
Define the adaptive small-angle scale
\[
\tau_{\theta,k}
=
\max\left\{
\tau_{\theta,\min},
\eta_\theta\,{\rm median}_{i\in\Edel}|\bar\theta_i|
\right\}.
\]
Here \(\tau_{\theta,k}\) is the step-dependent reference scale for deciding whether an angle is small, \(\tau_{\theta,\min}\) is the minimum allowed scale, \(\eta_\theta\) is a multiplier on the current median eligible angle, and the median is taken over eligible prune coordinates \(i\in\Edel\).
Then
\[
\boxed{
w_j^{\rm small}
=
\left[
1+
\left(
\frac{a_j^\theta}{\tau_{\theta,k}+\epsilon_\theta}
\right)^{p_\theta}
\right]^{-1}.
}
\]
Here \(w_j^{\rm small}\in(0,1]\) is large when \(|\bar\theta_j|\) is small compared with \(\tau_{\theta,k}\), \(\epsilon_\theta>0\) prevents division by zero, and \(p_\theta\) controls how sharply the weight decays for larger angles.
For age \(A_j(k)=k-a_j\),
\[
\boxed{
w_j^{\rm stale}
=
1+\eta_{\rm stale}
\left[
1-\exp\left(
-\frac{(A_j(k)-p_{\rm protect})_+}{\tau_{\rm age}}
\right)
\right].
}
\]
Here \(A_j(k)\) is the branch-local age of coordinate \(j\), \(w_j^{\rm stale}\) is the stale-coordinate preference, \(\eta_{\rm stale}\) sets the maximum stale-age boost, and \(\tau_{\rm age}\) controls how quickly the boost saturates after the protection window.

\blk{Physical interpretation}
Small and old coordinates become easier to nominate, but neither feature can bypass Schur loss, novelty, history, conditioning, or measured delete safety.

\section*{6. Deletion Novelty as Residual Tangent Fraction}

\blk{Micro-intuition}
For pruning, novelty reverses meaning: high novelty protects a coordinate because survivors cannot reproduce its tangent direction.

\blk{Mathematical derivation}
Let \(\tilde t_i\) be the horizontal Fubini--Study tangent associated with coordinate \(d_i\). For survivor set \(W_j\),
\[
G_{jj}=\operatorname{Re}\langle \tilde t_j,\tilde t_j\rangle,
\qquad
G_{jW}=
\bigl(\operatorname{Re}\langle \tilde t_j,\tilde t_i\rangle\bigr)_{i\in W_j},
\]
\[
G_{WW}=
\bigl(\operatorname{Re}\langle \tilde t_i,\tilde t_\ell\rangle\bigr)_{i,\ell\in W_j}.
\]
Here \(W_j\) is the chosen survivor coordinate window used to ask whether the remaining ansatz directions can reproduce the tangent of \(d_j\). \(G_{jj}\) is the Fubini--Study norm of the deleted coordinate's tangent, \(G_{jW}\) is the row vector of overlaps between the deleted tangent and survivor tangents, \(G_{Wj}=G_{jW}^{\top}\), and \(G_{WW}\) is the survivor Gram block.
Survivors explain \(d_j\)'s tangent by the metric-corrected projection fraction
\[
\eta_j^{\rm span}
=
\frac{
G_{jW}(G_{WW}+\lambda_G I)^+G_{Wj}
}{
G_{jj}+\epsilon_G
}.
\]
Here \(\eta_j^{\rm span}\) is the fraction of the deleted tangent explained by the survivor window, \(\lambda_G\) is the Gram ridge used in the pseudoinverse solve, \(I\) is the identity on the survivor window, \((\cdot)^+\) is the Moore--Penrose pseudoinverse, and \(\epsilon_G>0\) stabilizes the normalization.
The deletion novelty is the residual fraction
\[
\boxed{
\nudel
=
\clip\left(1-\eta_j^{\rm span}\right).
}
\]
Here \(\clip(x)=\min\{1,\max\{0,x\}\}\) clips the residual into \([0,1]\), and \(\nudel\) is the unexplained tangent fraction left after the survivor projection.
Then a smooth redundancy preference is
\[
\boxed{
w_j^{\rm red}
=
\sigma\left(
\frac{\nu_\star-\nudel}{\tau_\nu}
\right),
\qquad
\sigma(x)=\frac{1}{1+e^{-x}}.
}
\]
Here \(w_j^{\rm red}\in(0,1)\) is large when deletion novelty is below the redundancy target \(\nu_\star\), \(\tau_\nu\) sets the softness of the transition, and \(\sigma\) is the logistic function.

\blk{Physical interpretation}
\(\nudel\approx0\) means the survivors span the coordinate's tangent well, so deletion is plausible. \(\nudel\approx1\) means deletion removes a structurally distinct tangent direction, so the coordinate should be protected.

\section*{7. Persistence Gate}

\blk{Micro-intuition}
Persistence should stabilize prune nominations without requiring the exact same full batch to recur.

\blk{Mathematical derivation}
Let \(\mathcal N_r^\ominus\) be the scored prune-nomination list at step \(r\), before measured delete-refit. Let \(\phi(j)\) be a family key for coordinate identity, label, Pauli support family, or operator class. Define
\[
\boxed{
p_j(k)
=
\sum_{r=\max\{0,k-W_p\}}^{k-1}
\mathbf 1\left\{
j\in\mathcal N_r^\ominus
\;\lor\;
\phi(j)\in\phi(\mathcal N_r^\ominus)
\right\}.
}
\]
Here \(p_j(k)\) counts how often coordinate \(j\), or a coordinate in the same family \(\phi(j)\), appeared in recent prune-nomination lists. \(W_p\) is the persistence lookback window, \(r\) indexes prior adaptive steps, \(\mathcal N_r^\ominus\) is the prior scored nomination list, and \(\phi(\mathcal N_r^\ominus)\) is the set of family keys represented in that list.
Measured deletion is allowed only when
\[
\boxed{
p_j(k)\ge p_{\min}.
}
\]
Here \(p_{\min}\) is the required number of recent nomination appearances before the expensive measured delete-refit trial is allowed.
Thus the selection sequence is
\[
\begin{aligned}
j_k^\star\in\arg\max_{j\in\Gdel}\Rdel,
\\
p_{j_k^\star}(k)<p_{\min}
\;&\Rightarrow\;
\text{no measured delete call}.
\end{aligned}
\]
Here the first line chooses the highest-scoring candidate, and the second line suppresses the measured delete test if that candidate has not persisted long enough.

\blk{Physical interpretation}
The scalar score still chooses the candidate. Persistence decides whether the score has been stable enough to justify paying the measured delete-refit cost.

\section*{8. Rejection History Penalty}

\blk{Micro-intuition}
Cooldown blocks immediate retries; rejection history should make repeatedly unsafe deletions harder to nominate later.

\blk{Mathematical derivation}
After a rejected measured deletion trial for \(j\),
\[
\boxed{
\bar L_{j,k+1}^{\rm rej}
=
(1-\rho_{\rm rej})\bar L_{j,k}^{\rm rej}
+\rho_{\rm rej}
\max\{0,\dEdel-\epsilon_k^\ominus\}.
}
\]
Here \(\bar L_{j,k}^{\rm rej}\) is the accumulated rejection-loss history for coordinate \(j\) before step \(k\), \(\bar L_{j,k+1}^{\rm rej}\) is its updated value, and \(\rho_{\rm rej}\in[0,1]\) is the update rate. The term \(\max\{0,\dEdel-\epsilon_k^\ominus\}\) records only the amount by which the rejected deletion exceeded the allowed prune tolerance.
If \(j\) is not measured at step \(k\), use idle decay
\[
\bar L_{j,k+1}^{\rm rej}
=
(1-\rho_{\rm idle})\bar L_{j,k}^{\rm rej},
\qquad
0\le\rho_{\rm idle}\le\rho_{\rm rej}.
\]
Here \(\rho_{\rm idle}\) is the decay rate for rejection history on steps where the coordinate is not tested.
This penalty enters \(R_j^\ominus\) through
\[
\epsilon_R+\cdots+\lambda_{\rm rej}\bar L_j^{\rm rej}+\cdots.
\]
Here the ellipses denote the other denominator terms already defined in the scalar score.

\blk{Physical interpretation}
A badly rejected prune candidate remains possible later, but it must overcome its own rejection history. A nearly safe rejection receives only a small future penalty.

\section*{9. Conditioning Pressure}

\blk{Micro-intuition}
Deleting a coordinate can improve or damage tangent-space conditioning; this should first be telemetry, then optionally a risk term.

\blk{Mathematical derivation}
Let \(G_k\) be the Fubini--Study Gram matrix for the active ansatz tangents and \(G_{k,-j}\) the survivor Gram matrix after removing coordinate \(j\). With ridge \(\lambda_\kappa\),
\[
G_{k,\lambda}=G_k+\lambda_\kappa I,
\qquad
G_{k,-j,\lambda}=G_{k,-j}+\lambda_\kappa I.
\]
Here \(G_{k,\lambda}\) and \(G_{k,-j,\lambda}\) are ridge-regularized Gram matrices before and after deleting \(d_j\), and \(\lambda_\kappa\) is the conditioning ridge.
Define condition-number damage
\[
\boxed{
\Delta\kappa_j
=
\left[
\log\kappa(G_{k,-j,\lambda})
-
\log\kappa(G_{k,\lambda})
\right]_+.
}
\]
Here \(\kappa(\cdot)\) is the matrix condition number, \([\cdot]_+=\max\{0,\cdot\}\), and \(\Delta\kappa_j\) measures only worsening of conditioning after deletion.
Define rank-loss telemetry
\[
\boxed{
\Delta r_j
=
{\rm rank}_{\tau}(G_k)
-
{\rm rank}_{\tau}(G_{k,-j}).
}
\]
Here \({\rm rank}_{\tau}\) is numerical rank at threshold \(\tau\), and \(\Delta r_j\) counts how many numerically independent tangent directions would be lost by deleting \(d_j\).

\blk{Physical interpretation}
\(\Delta\kappa_j=0\) means deletion does not worsen conditioning by this metric. \(\Delta r_j>0\) means deletion removes a numerically independent tangent direction, which overlaps with high deletion novelty.

\section*{10. Schur Accuracy Knobs}

\blk{Micro-intuition}
The scalar prune score depends on \(\Lam\), so Schur prediction quality must be audited.

\blk{Mathematical derivation}
Recommended Schur-window sweep:
\[
w_{\rm Schur}\in\{4,8,12,{\rm full}\}.
\]
Here \(w_{\rm Schur}\) is the number of survivor coordinates allowed in the local Schur model; \({\rm full}\) means all surviving coordinates may enter the Schur block.
For even \(w\), a centered interior window is
\[
W_j^{\rm centered}
=
\{d_{j-w/2},\ldots,d_{j-1},d_{j+1},\ldots,d_{j+w/2}\},
\]
with boundary shifts near the ends.
Here \(w\) is a chosen finite Schur-window size, \(W_j^{\rm centered}\) is the centered survivor window around deletion coordinate \(j\), and \(d_j\) itself is omitted because it is the coordinate being tested for deletion.

Recommended measured trial cap sweep:
\[
M_{\rm exact}^{\ominus}\in\{1,2,3\}.
\]
Here \(M_{\rm exact}^{\ominus}\) is the maximum number of Schur-nominated prune candidates that may be sent to measured delete-refit at one adaptive step.
For every measured prune trial, record
\[
\boxed{
(\Lam,\dEdel,\dEdel-\Lam,A_j^\ominus).
}
\]
Here the recorded tuple compares the Schur predicted loss, the measured deletion loss, the prediction error, and the final accept/reject result for the same coordinate \(j\).

\blk{Physical interpretation}
The dangerous failure mode is a low predicted \(\Lam\) with high measured \(\dEdel\). Window size, centering, and measured cap should be tuned by that diagnostic, not by final energy alone.

\section*{11. Measurement-Cost Boundary}

\blk{Micro-intuition}
Energy regression is already the measured prune safety test; extra overlap tests should not silently become hardware-facing acceptance guards.

\blk{Mathematical derivation}
Default acceptance remains
\[
\dEdel
=
E(\Ok\ominus d_j,\boldsymbol\theta_{k,-j})-\Ek
\le
\epsilon_k^\ominus
\]
together with retained-gain and curvature guards. Geometry-only quantities such as \(\nudel\) and \(\Delta\kappa_j\) are measurement-free only when their Gram blocks are already estimated or cached:
Here the inequality is the energy-regression guard: deleting \(d_j\) is allowed only if the measured energy increase stays below the prune tolerance \(\epsilon_k^\ominus\).
\[
\begin{gathered}
G_{jj},G_{jW},G_{WW},G_k,G_{k,-j}\\
\subseteq
\text{available SNAKE geometry cache}.
\end{gathered}
\]
Here the listed Gram objects are the tangent-overlap quantities needed for novelty and conditioning telemetry. The subset relation means these quantities are already available from the algorithm's geometry cache; if not, estimating them introduces additional measurement cost.
Otherwise, they require additional metric/overlap estimation.

\blk{Physical interpretation}
Use state-overlap and fidelity checks as simulator diagnostics by default. Make them hardware-facing acceptance guards only after an explicit measurement-budget argument.

\section*{12. Proposed Prune Flow}

\blk{Micro-intuition}
The redesign changes nomination, not the measured deletion authority.

\blk{Mathematical derivation}
\[
\boxed{
\begin{gathered}
\Edel
\xrightarrow{\rm features}
\Gdel
\xrightarrow{\max R_j^\ominus}
j_k^\star
\\
\xrightarrow{p_j(k)\ge p_{\min}}
\Cdel
\xrightarrow{\rm measured}
A_j^\ominus.
\end{gathered}
}
\]
Here the flow reads from left to right: eligible coordinates are featurized, hard gates produce the scoreable set, the scalar score chooses a top nomination, persistence decides whether the nomination is stable enough to measure, and the measured delete-refit test produces \(A_j^\ominus\).
Equivalently,
\[
\boxed{
\begin{gathered}
j\in\Edel
\xrightarrow{\rm score}
R_j^\ominus
\xrightarrow{\rm persistence}
\Cdel
\xrightarrow{p_j(k)\ge p_{\min}}
A_j^\ominus.
\end{gathered}
}
\]
Here the second display emphasizes the per-coordinate view: an eligible coordinate \(j\) receives a score \(R_j^\ominus\), may enter the measured candidate set \(\Cdel\), and is finally accepted or rejected by \(A_j^\ominus\).
On rejection,
\[
\begin{aligned}
A_j^\ominus=0
\quad&\Longrightarrow\quad
\Ok^{\rm out}=\Ok,
\\
&c_j(k+1)=c_{\rm cool},
\qquad
\bar L_{j,k+1}^{\rm rej}\ {\rm increases}.
\end{aligned}
\]
Here \(\Ok^{\rm out}\) is the ansatz returned by prune after the failed deletion trial, \(c_{\rm cool}\) is the cooldown reset value, and the rejection history increases according to the update in Sec.~8.
On acceptance,
\[
A_j^\ominus=1
\quad\Longrightarrow\quad
\Ok^{\rm out}=\Ok\ominus d_j.
\]
Here the accepted prune commits the deleted ansatz as the outgoing branch state.

\blk{Physical interpretation}
Nomination should become resource-aware, redundancy-aware, history-aware, and geometry-aware. Acceptance remains the measured delete-refit safety test.

\section*{13. Implementation-Later Hooks}

\blk{Micro-intuition}
The likely code changes are localized, but this PDF does not implement them.

\blk{Mathematical derivation}
Future hooks should map the mathematical objects as follows:
{\scriptsize
\[
\begin{array}{c|p{0.62\columnwidth}}
\text{object} & \text{likely implementation hook}\\
\hline
\Lam & \texttt{build\_static\_prune\_surrogate\_scores}\\
\Rdel & \texttt{rank\_prune\_candidates} replacement or extension\\
A_j^\ominus & \texttt{apply\_pruning}\text{ remains authoritative}\\
c_j,\bar L_j^{\rm rej} & prune metadata transport in \texttt{adapt\_pipeline.py}\\
\nu_j^\ominus,\Delta\kappa_j & \text{geometry/cache telemetry}\\
p_j(k) & \text{nomination-history state}
\end{array}
\]
}
Here the left column lists mathematical objects from this design note, and the right column lists likely implementation locations for a later code pass. This table is only a planning map; it does not change the current implementation.

\blk{Physical interpretation}
The first code pass should emit telemetry for every score component before using all gates aggressively. The benchmark contract should be rerun only after the component audit is interpretable.

\end{document}
